\section*{Conclusion générale et perspectives}
\addcontentsline{toc}{section}{Conclusion générale}

Ce projet de fin d'études a eu pour ambition la refonte architecturale et le déploiement automatisé d'un système de gestion et de facturation télécom (BSS). Face à la rigidité des systèmes monolithiques traditionnels, la solution proposée a permis de démontrer les bénéfices liés à l'adoption d'une architecture orientée microservices, combinée aux pratiques DevOps modernes.

Tout au long de ce projet, nous avons suivi une démarche agile et itérative, structurée en quatre releases majeures, garantissant une évolution maîtrisée du produit :
\begin{itemize}
    \item \textbf{La conception logicielle :} Nous avons scindé les domaines métier en microservices distincts (Customer, Boutique, Catalog, Subscription, Billing et Auth), assurant ainsi l'indépendance du code et des bases de données.
    \item \textbf{L'orchestration et le Cloud :} Le déploiement s'est concrétisé par la containerisation Docker et la mise en place d'une infrastructure robuste sur Amazon EKS, illustrant la maîtrise des écosystèmes Cloud Native.
    \item \textbf{Culture DevOps et automatisation :} La conception d'un pipeline CI/CD complet via Jenkins, intégrant les tests de sécurité (Trivy), d'analyse de code (SonarQube) et le paradigme GitOps (ArgoCD), a démontré notre capacité à automatiser la chaîne de livraison logicielle.
    \item \textbf{Supervision :} Enfin, la mise en place de Prometheus et Grafana a clôturé le cycle de vie du projet en fournissant la visibilité nécessaire pour garantir la stabilité de l'application en production.
\end{itemize}

Ce projet a été riche en défis techniques et en apprentissages. Il nous a permis d'approfondir nos compétences en ingénierie logicielle (Java, Spring Boot, Angular), en architecture distribuée et en administration d'infrastructures cloud et Kubernetes.

\subsection*{Limites et Perspectives}
Bien que la solution réponde aux objectifs initiaux de ce PoC (Proof of Concept), une prise de recul architecturale permet d'identifier certaines limites, représentant autant de perspectives d'amélioration pour une mise en production industrielle :

\begin{enumerate}
    \item \textbf{Couplage temporel (Le "Monolithe Distribué") :} Actuellement, la communication inter-services repose sur des appels synchrones via OpenFeign (HTTP/REST). Ce choix pragmatique crée un couplage fort : si un service échoue, l'appel en cascade échoue également. La solution radicale pour l'avenir sera la transition vers une \textbf{Architecture Orientée Événements (EDA)} utilisant un broker de messages (Apache Kafka), garantissant l'asynchronisme et la haute résilience.
    
    \item \textbf{Cohérence des données (Saga Pattern) :} En l'absence de transactions de bases de données globales, les erreurs partielles peuvent entraîner des incohérences. L'implémentation du patron \textbf{Saga} permettrait d'orchestrer des transactions distribuées via des requêtes de compensation.
    
    \item \textbf{Gestion des performances et latence réseau :} Les multiples appels réseau nécessaires pour la facturation ralentissent le système. La mise en place du pattern \textbf{CQRS} (Command Query Responsibility Segregation) et d'un système de \textbf{Cache distribué} (Redis) pour des données peu volatiles (comme le catalogue d'offres) réduirait considérablement la latence, à condition de maîtriser l'invalidation de cache pour éviter toute erreur de tarification.
\end{enumerate}

En définitive, ce projet de fin d'études constitue une fondation technique extrêmement solide, alignée sur les standards actuels de l'industrie informatique. Les perspectives identifiées prouvent que l'ingénierie logicielle est un processus d'amélioration continue, visant sans cesse à allier résilience opérationnelle et excellence technique.

\newpage
